\documentclass[a4paper,12pt]{article}
\usepackage{amsmath, amssymb}
\usepackage{xcolor}
\usepackage[most]{tcolorbox}
\usepackage{geometry}
\graphicspath{ {../Images/} }

\geometry{left=1.5cm, right=1.5cm, top=1.5cm, bottom=1.5cm}

\title{Exercices de Polynômes de Degré 2 - Niveau Première}
\author{}
\date{}

% Style de boîte pour les exercices
\tcbset{
    colframe=purple!70,
    colback=white,
    coltitle=black,
    fonttitle=\bfseries,
    sharp corners,
    boxrule=0.5mm,
    colbacktitle=purple!20!white,
    center title,
    boxsep=4pt,
}

% Style de boîte pour les solutions
\newtcolorbox{solutionbox}[1][]{
    colframe=green!70,
    colback=white,
    coltitle=black,
    fonttitle=\bfseries,
    sharp corners,
    boxrule=0.5mm,
    colbacktitle=green!20!white,
    title=Solution,
    boxsep=4pt,
    breakable,
    #1
}

\begin{document}

% Exercice 1
\begin{tcolorbox}[title=Exercice 1 * {Calculer}]
Calculer le discriminant de chaque polynôme ci-dessous.
\begin{enumerate}
    \item $3x^2 + 2x + 1$
    \item $2x^2 - 2x - 4$
    \item $-3x^2 + 2x - 5$
    \item $x^2 + x + 1$
    \item $x^2 - 7$
    \item $x - x^2 + 3$
    \item $7x - 3 + x^2$
    \item $1 - 3x^2$
\end{enumerate}
\end{tcolorbox}

\begin{solutionbox}
Pour chaque polynôme de la forme $ax^2 + bx + c$, le discriminant est $\Delta = b^2 - 4ac$.

\begin{enumerate}
    \item $3x^2 + 2x + 1$

    Coefficients : $a = 3$, $b = 2$, $c = 1$

    Calcul du discriminant :
\[
    \Delta = b^2 - 4ac = (2)^2 - 4 \times 3 \times 1 = 4 - 12 = -8
\]

    \item $2x^2 - 2x - 4$

    Coefficients : $a = 2$, $b = -2$, $c = -4$

    Calcul du discriminant :
\[
    \Delta = (-2)^2 - 4 \times 2 \times (-4) = 4 + 32 = 36
\]

    \item $-3x^2 + 2x - 5$

    Coefficients : $a = -3$, $b = 2$, $c = -5$

    Calcul du discriminant :
\[
    \Delta = (2)^2 - 4 \times (-3) \times (-5) = 4 - 60 = -56
\]

    \item $x^2 + x + 1$

    Coefficients : $a = 1$, $b = 1$, $c = 1$

    Calcul du discriminant :
\[
    \Delta = (1)^2 - 4 \times 1 \times 1 = 1 - 4 = -3
\]

    \item $x^2 - 7$

    Coefficients : $a = 1$, $b = 0$, $c = -7$

    Calcul du discriminant :
\[
    \Delta = (0)^2 - 4 \times 1 \times (-7) = 0 + 28 = 28
\]

    \item $x - x^2 + 3$

    Réécrivons le polynôme sous la forme standard :
\[
    -x^2 + x + 3
\]
    Coefficients : $a = -1$, $b = 1$, $c = 3$

    Calcul du discriminant :
\[
    \Delta = (1)^2 - 4 \times (-1) \times 3 = 1 + 12 = 13
\]

    \item $7x - 3 + x^2$

    Réécrivons le polynôme :
\[
    x^2 + 7x - 3
\]
    Coefficients : $a = 1$, $b = 7$, $c = -3$

    Calcul du discriminant :
\[
    \Delta = (7)^2 - 4 \times 1 \times (-3) = 49 + 12 = 61
\]

    \item $1 - 3x^2$

    Réécrivons le polynôme :
\[
    -3x^2 + 1
\]
    Coefficients : $a = -3$, $b = 0$, $c = 1$

    Calcul du discriminant :
\[
    \Delta = (0)^2 - 4 \times (-3) \times 1 = 0 + 12 = 12
\]
\end{enumerate}
\end{solutionbox}

% Exercice 2
\begin{tcolorbox}[title=Exercice 2 * {Calculer}]
Déterminer la forme factorisée, lorsqu’elle existe, des polynômes suivants.
\begin{enumerate}
    \item $x^2 + 3x - 4$
    \item $2x^2 - 4x - 12$
    \item $-3x^2 + x - 5$
    \item $x^2 + x + 1$
    \item $x^2 - 2x + 1$
    \item $x - 4x^2 + 2$
\end{enumerate}
\end{tcolorbox}

\begin{solutionbox}
Pour factoriser un polynôme de degré 2, on calcule son discriminant $\Delta$.

\begin{enumerate}
    \item $x^2 + 3x - 4$

    $\Delta = 3^2 - 4 \times 1 \times (-4) = 9 + 16 = 25$

    Les racines sont :
\[
    x = \dfrac{-b \pm \sqrt{\Delta}}{2a} = \dfrac{-3 \pm 5}{2}
\]
    Donc :
\[
    x_1 = \dfrac{-3 + 5}{2} = 1 \quad \text{et} \quad x_2 = \dfrac{-3 - 5}{2} = -4
\]

    La forme factorisée est :
\[
    x^2 + 3x - 4 = (x - x_1)(x - x_2) = (x - 1)(x + 4)
\]

    \item $2x^2 - 4x - 12$

    $\Delta = (-4)^2 - 4 \times 2 \times (-12) = 16 + 96 = 112$

    Les racines sont :
\[
    x = \dfrac{4 \pm \sqrt{112}}{4} = \dfrac{4 \pm 4 \sqrt{7}}{4} = \dfrac{4(1 \pm \sqrt{7})}{4} = 1 \pm \sqrt{7}
\]

    La forme factorisée est :
\[
    2x^2 - 4x - 12 = 2 \left( x - \left( 1 + \sqrt{7} \right) \right) \left( x - \left( 1 - \sqrt{7} \right) \right)
\]

    \item $-3x^2 + x - 5$

    $\Delta = 1^2 - 4 \times (-3) \times (-5) = 1 - 60 = -59$

    Comme $\Delta < 0$, le polynôme n'est pas factorisable dans $\mathbb{R}$.

    \item $x^2 + x + 1$

    $\Delta = 1^2 - 4 \times 1 \times 1 = 1 - 4 = -3$

    Comme $\Delta < 0$, le polynôme n'est pas factorisable dans $\mathbb{R}$.

    \item $x^2 - 2x + 1$

    $\Delta = (-2)^2 - 4 \times 1 \times 1 = 4 - 4 = 0$

    Une racine double :
\[
    x = \dfrac{2}{2} = 1
\]

    La forme factorisée est :
\[
    x^2 - 2x + 1 = (x - 1)^2
\]

    \item $x - 4x^2 + 2$

    Réécrivons le polynôme :
\[
    -4x^2 + x + 2
\]

    Coefficients : $a = -4$, $b = 1$, $c = 2$

    $\Delta = 1^2 - 4 \times (-4) \times 2 = 1 + 32 = 33$

    Racines :
\[
    x = \dfrac{-1 \pm \sqrt{33}}{-8} = \dfrac{-1 \pm \sqrt{33}}{-8}
\]

    La forme factorisée est :
\[
    -4x^2 + x + 2 = -4 \left( x - x_1 \right) \left( x - x_2 \right)
\]

    avec $x_1 = \dfrac{-1 + \sqrt{33}}{-8}$ et $x_2 = \dfrac{-1 - \sqrt{33}}{-8}$.

\end{enumerate}
\end{solutionbox}

% Exercice 3
\begin{tcolorbox}[title=Exercice 3 * {Calculer}]
Résoudre dans $ \mathbb{R} $ les équations suivantes.
\begin{enumerate}
    \item $x^2 + x - 4 = 0$
    \item $3x^2 - x - 6 = 0$
    \item $x^2 - 10x + 25 = 0$
    \item $x^2 - 3x + 3 = 0$
    \item $x^2 - 1 = 0$
    \item $x - 4x^2 + 7 = 0$
\end{enumerate}
\end{tcolorbox}

\begin{solutionbox}
Résolvons chaque équation en calculant le discriminant.

\begin{enumerate}
    \item $x^2 + x - 4 = 0$

    $\Delta = 1^2 - 4 \times 1 \times (-4) = 1 + 16 = 17$

    Racines :
\[
    x = \dfrac{-1 \pm \sqrt{17}}{2}
\]

    \item $3x^2 - x - 6 = 0$

    $\Delta = (-1)^2 - 4 \times 3 \times (-6) = 1 + 72 = 73$

    Racines :
\[
    x = \dfrac{1 \pm \sqrt{73}}{6}
\]

    \item $x^2 - 10x + 25 = 0$

    $\Delta = (-10)^2 - 4 \times 1 \times 25 = 100 - 100 = 0$

    Racine double :
\[
    x = \dfrac{10}{2} = 5
\]

    \item $x^2 - 3x + 3 = 0$

    $\Delta = (-3)^2 - 4 \times 1 \times 3 = 9 - 12 = -3$

    Pas de solution réelle.

    \item $x^2 - 1 = 0$

    $\Delta = 0^2 - 4 \times 1 \times (-1) = 0 + 4 = 4$

    Racines :
\[
    x = \dfrac{0 \pm 2}{2} = \pm 1
\]

    \item $x - 4x^2 + 7 = 0$

    Réécrivons l'équation :
\[
    -4x^2 + x + 7 = 0
\]

    $\Delta = 1^2 - 4 \times (-4) \times 7 = 1 + 112 = 113$

    Racines :
\[
    x = \dfrac{-1 \pm \sqrt{113}}{-8}
\]
\end{enumerate}
\end{solutionbox}

% Exercice 4
\begin{tcolorbox}[title=Exercice 4 * {Calculer}]
Déterminer les racines des polynômes suivants.
\begin{enumerate}
    \item $x^2 - x - 1$
    \item $x^2 - x$
    \item $-5x^2 + 7x + 12$
    \item $-3x - x^2 + 4$
    \item $3x^2 - 6x + 3$
    \item $-x + x^2 + 7$
\end{enumerate}
\end{tcolorbox}

\begin{solutionbox}
\begin{enumerate}
    \item $x^2 - x - 1$

    $\Delta = (-1)^2 - 4 \times 1 \times (-1) = 1 + 4 = 5$

    Racines :
\[
    x = \dfrac{1 \pm \sqrt{5}}{2}
\]

    \item $x^2 - x$

    Mise en évidence :
\[
    x(x - 1) = 0
\]
    Les racines sont $x = 0$ et $x = 1$.

    \item $-5x^2 + 7x + 12$

    $\Delta = 7^2 - 4 \times (-5) \times 12 = 49 + 240 = 289$

    Racines :
\[
    x = \dfrac{-7 \pm 17}{-10}
\]

    Calculons :
\[
    x_1 = \dfrac{-7 + 17}{-10} = \dfrac{10}{-10} = -1
\]
\[
    x_2 = \dfrac{-7 - 17}{-10} = \dfrac{-24}{-10} = \dfrac{12}{5}
\]

    \item $-3x - x^2 + 4$

    Réécrivons :
\[
    -x^2 - 3x + 4 = 0
\]
    Puis multiplions par $-1$ :
\[
    x^2 + 3x - 4 = 0
\]

    $\Delta = 3^2 - 4 \times 1 \times (-4) = 9 + 16 = 25$

    Racines :
\[
    x = \dfrac{-3 \pm 5}{2}
\]

    Donc :
\[
    x_1 = \dfrac{-3 + 5}{2} = 1 \quad \text{et} \quad x_2 = \dfrac{-3 - 5}{2} = -4
\]

    \item $3x^2 - 6x + 3$

    $\Delta = (-6)^2 - 4 \times 3 \times 3 = 36 - 36 = 0$

    Racine double :
\[
    x = \dfrac{6}{6} = 1
\]

    \item $-x + x^2 + 7$

    Réécrivons :
\[
    x^2 - x + 7 = 0
\]

    $\Delta = (-1)^2 - 4 \times 1 \times 7 = 1 - 28 = -27$

    Pas de racines réelles.

\end{enumerate}
\end{solutionbox}

% Exercice 5
\begin{tcolorbox}[title=Exercice 5 ** {Calculer, Raisonner}]
Soit $m \in \mathbb{R}$. On considère l'équation $x^2 - mx + m + 2 = 0$. Pour quelles valeurs de $m$ l'équation admet-elle une unique solution ?
\end{tcolorbox}

\begin{solutionbox}
Une équation du second degré admet une unique solution réelle lorsque son discriminant est nul.

Calculons le discriminant :
\[
\Delta = (-m)^2 - 4 \times 1 \times (m + 2) = m^2 - 4(m + 2)
\]

Développons :
\[
\Delta = m^2 - 4m - 8
\]

Posons $\Delta = 0$ :
\[
m^2 - 4m - 8 = 0
\]

Résolvons cette équation du second degré en $m$ :
\[
\Delta' = (-4)^2 - 4 \times 1 \times (-8) = 16 + 32 = 48
\]

Racines :
\[
m = \dfrac{4 \pm \sqrt{48}}{2} = \dfrac{4 \pm 4 \sqrt{3}}{2} = 2 \pm 2 \sqrt{3}
\]

Ainsi, l'équation admet une unique solution réelle pour :
\[
m = 2 + 2 \sqrt{3} \quad \text{ou} \quad m = 2 - 2 \sqrt{3}
\]
\end{solutionbox}

% Exercice 6
\begin{tcolorbox}[title=Exercice 6 ** {Calculer, Modéliser}]
Déterminer deux entiers naturels consécutifs dont le produit est égal à 64262.
\end{tcolorbox}

\begin{solutionbox}
Soit $n$ un entier naturel. Les deux entiers consécutifs sont $n$ et $n + 1$.

Leur produit est :
\[
n(n + 1) = 64262
\]

Développons :
\[
n^2 + n - 64262 = 0
\]

Résolvons cette équation du second degré :

$\Delta = 1^2 - 4 \times 1 \times (-64262) = 1 + 257048 = 257049$

$\sqrt{\Delta} = \sqrt{257049} = 507$

Racines :
\[
n = \dfrac{-1 \pm 507}{2}
\]

Calculons :

1) $n = \dfrac{-1 + 507}{2} = \dfrac{506}{2} = 253$

2) $n = \dfrac{-1 - 507}{2} = \dfrac{-508}{2} = -254$ (négatif, pas un entier naturel)

Donc, les deux entiers sont $253$ et $254$.

\textbf{Réponse :} Les entiers sont $253$ et $254$.
\end{solutionbox}

% Exercice 7
\begin{tcolorbox}[title=Exercice 7 * {Calculer}]
Les équations suivantes admettent toutes des racines. Dans chaque cas, déterminer le produit et la somme des racines.
\begin{enumerate}
    \item $x^2 - x - 2$
    \item $x^2 - \pi x$
    \item $-x^2 + 3x + 12$
    \item $3x + x^2 - 1$
    \item $\sqrt{17}x^2 + 32x - 1$
    \item $2x - 3x^2 + 5$
\end{enumerate}
\end{tcolorbox}

\begin{solutionbox}
Pour un polynôme $ax^2 + bx + c$, la somme des racines est $S = -\dfrac{b}{a}$ et le produit est $P = \dfrac{c}{a}$.

\begin{enumerate}
    \item $x^2 - x - 2$

    Coefficients : $a = 1$, $b = -1$, $c = -2$

    Somme :
\[
    S = -\dfrac{-1}{1} = 1
\]

    Produit :
\[
    P = \dfrac{-2}{1} = -2
\]

    \item $x^2 - \pi x$

    Coefficients : $a = 1$, $b = -\pi$, $c = 0$

    Somme :
\[
    S = -\dfrac{-\pi}{1} = \pi
\]

    Produit :
\[
    P = \dfrac{0}{1} = 0
\]

    \item $-x^2 + 3x + 12$

    Coefficients : $a = -1$, $b = 3$, $c = 12$

    Somme :
\[
    S = -\dfrac{3}{-1} = 3
\]

    Produit :
\[
    P = \dfrac{12}{-1} = -12
\]

    \item $3x + x^2 - 1$

    Réécrivons :
\[
    x^2 + 3x - 1
\]
    Coefficients : $a = 1$, $b = 3$, $c = -1$

    Somme :
\[
    S = -\dfrac{3}{1} = -3
\]

    Produit :
\[
    P = \dfrac{-1}{1} = -1
\]

    \item $\sqrt{17}x^2 + 32x - 1$

    Coefficients : $a = \sqrt{17}$, $b = 32$, $c = -1$

    Somme :
\[
    S = -\dfrac{32}{\sqrt{17}}
\]

    Produit :
\[
    P = \dfrac{-1}{\sqrt{17}}
\]

    \item $2x - 3x^2 + 5$

    Réécrivons :
\[
    -3x^2 + 2x + 5
\]
    Coefficients : $a = -3$, $b = 2$, $c = 5$

    Somme :
\[
    S = -\dfrac{2}{-3} = \dfrac{2}{3}
\]

    Produit :
\[
    P = \dfrac{5}{-3} = -\dfrac{5}{3}
\]
\end{enumerate}
\end{solutionbox}

% Exercice 8
\begin{tcolorbox}[title=Exercice 8 ** {Calculer}]
Déterminer deux réels $x_1$ et $x_2$ vérifiant les conditions suivantes :
\begin{equation*}
\begin{cases}
x_1 + x_2 = 5 \\ 
x_1 x_2 = 1
\end{cases}
\end{equation*}
\end{tcolorbox}

\begin{solutionbox}
Les nombres $x_1$ et $x_2$ sont les racines du polynôme :
\[
x^2 - (x_1 + x_2) x + x_1 x_2 = 0
\]
Substituons les valeurs :
\[
x^2 - 5x + 1 = 0
\]
Résolvons cette équation :
\[
\Delta = (-5)^2 - 4 \times 1 \times 1 = 25 - 4 = 21
\]
Racines :
\[
x = \dfrac{5 \pm \sqrt{21}}{2}
\]
Donc :
\[
x_1 = \dfrac{5 + \sqrt{21}}{2}, \quad x_2 = \dfrac{5 - \sqrt{21}}{2}
\]
\end{solutionbox}

% Exercice 9
\begin{tcolorbox}[title=Exercice 9 ** {Calculer}]
Existe-t-il des réels $a$ et $b$ vérifiant les conditions suivantes ?
\begin{equation*}
\begin{cases}
a + b = -3 \\ 
ab = 2
\end{cases}
\end{equation*}
Si oui, déterminer lesquels.
\end{tcolorbox}

\begin{solutionbox}
Les nombres $a$ et $b$ sont les racines du polynôme :
\[
x^2 - (a + b) x + ab = 0
\]
Substituons les valeurs :
\[
x^2 - (-3)x + 2 = x^2 + 3x + 2 = 0
\]
Résolvons :
\[
\Delta = 3^2 - 4 \times 1 \times 2 = 9 - 8 = 1
\]
Racines :
\[
x = \dfrac{-3 \pm 1}{2}
\]
Calculons :
\[
x_1 = \dfrac{-3 + 1}{2} = -1
\]
\[
x_2 = \dfrac{-3 - 1}{2} = -2
\]
Donc, $a$ et $b$ sont $-1$ et $-2$ (dans n'importe quel ordre).

\textbf{Réponse :} Oui, $a = -1$, $b = -2$ ou $a = -2$, $b = -1$.
\end{solutionbox}

% Exercice 10
\begin{tcolorbox}[title=Exercice 10 ** {Calculer}]
Existe-t-il des réels $a$ et $b$ vérifiant les conditions suivantes ?
\begin{equation*}
\begin{cases}
a + b = 4 \\ 
ab = 5
\end{cases}
\end{tcolorbox}

\begin{solutionbox}
Les nombres $a$ et $b$ sont les racines du polynôme :
\[
x^2 - (a + b) x + ab = x^2 - 4x + 5 = 0
\]
Résolvons :
\[
\Delta = (-4)^2 - 4 \times 1 \times 5 = 16 - 20 = -4
\]
Comme $\Delta < 0$, il n'y a pas de racines réelles.

\textbf{Réponse :} Non, il n'existe pas de réels $a$ et $b$ vérifiant ces conditions.
\end{solutionbox}

% Exercice 11
\begin{tcolorbox}[title=Exercice 11 *** {Calculer, Chercher}]
Soient les réels $a$ et $b$ les solutions de l’équation $2x^2 - 7x + 5 = 0$. Calculer $a^2 + b^2$.
\end{tcolorbox}

\begin{solutionbox}
Les sommes et produit des racines sont :
\[
S = a + b = \dfrac{7}{2}, \quad P = ab = \dfrac{5}{2}
\]
Calculons $a^2 + b^2$ :
\[
a^2 + b^2 = (a + b)^2 - 2ab = S^2 - 2P
\]
Substituons les valeurs :
\[
a^2 + b^2 = \left( \dfrac{7}{2} \right )^2 - 2 \times \dfrac{5}{2} = \dfrac{49}{4} - 5 = \dfrac{49}{4} - \dfrac{20}{4} = \dfrac{29}{4}
\]
\textbf{Réponse :} $a^2 + b^2 = \dfrac{29}{4}$
\end{solutionbox}

% Exercice 12
\begin{tcolorbox}[title=Exercice 12 ** {Calculer}]
Résoudre les équations suivantes en utilisant la méthode la plus adaptée (facteur commun, identité remarquable, calcul avec le discriminant, racine évidente, etc.)
\begin{enumerate}
    \item $x^2 + 4x = 0$
    \item $\sqrt{2}x^2 + 3x - 1 = 0$
    \item $\pi x^2 + 2\pi x + \pi = 0$
    \item $x^2 + x + 1 = 0$
    \item $x^2 - x - 2 = 0$
    \item $x^2 + 7 = 0$
    \item $x^2 - 12x = -36$
    \item $-3x^2 - 1 = 0$
    \item $3x^2 + 3x + 0 = 0$
    \item $7x^2 + x + 2 = 0$
    \item $(x - 5)(x + 7) = 0$
    \item $(x - 3)(x + 1) = 6$
    \item $x^2 = \sqrt{2}$
\end{enumerate}
\end{tcolorbox}

\begin{solutionbox}
\begin{enumerate}
    \item $x^2 + 4x = 0$

    Mise en évidence :
\[
    x(x + 4) = 0
\]
    Solutions :
\[
    x = 0 \quad \text{ou} \quad x = -4
\]

    \item $\sqrt{2}x^2 + 3x - 1 = 0$

    Discriminant :
\[
    \Delta = 3^2 - 4 \times \sqrt{2} \times (-1) = 9 + 4\sqrt{2}
\]
    Racines :
\[
    x = \dfrac{-3 \pm \sqrt{9 + 4\sqrt{2}}}{2\sqrt{2}}
\]
    (Les racines sont irrationnelles.)

    \item $\pi x^2 + 2\pi x + \pi = 0$

    Divisons par $\pi$ :
\[
    x^2 + 2x + 1 = 0
\]
    C'est une identité remarquable :
\[
    (x + 1)^2 = 0
\]
    Solution double :
\[
    x = -1
\]

    \item $x^2 + x + 1 = 0$

    Discriminant :
\[
    \Delta = 1^2 - 4 \times 1 \times 1 = -3
\]
    Pas de solution réelle.

    \item $x^2 - x - 2 = 0$

    Discriminant :
\[
    \Delta = (-1)^2 - 4 \times 1 \times (-2) = 1 + 8 = 9
\]
    Racines :
\[
    x = \dfrac{1 \pm 3}{2}
\]
    Donc :
\[
    x = 2 \quad \text{ou} \quad x = -1
\]

    \item $x^2 + 7 = 0$

    Discriminant :
\[
    \Delta = 0^2 - 4 \times 1 \times 7 = -28
\]
    Pas de solution réelle.

    \item $x^2 - 12x = -36$

    Réécrivons :
\[
    x^2 - 12x + 36 = 0
\]
    Identité remarquable :
\[
    (x - 6)^2 = 0
\]
    Solution double :
\[
    x = 6
\]

    \item $-3x^2 - 1 = 0$

\[
    -3x^2 = 1 \quad \Rightarrow \quad x^2 = -\dfrac{1}{3}
\]
    Pas de solution réelle.

    \item $3x^2 + 3x + 0 = 0$

    Mise en évidence :
\[
    3x(x + 1) = 0
\]
    Solutions :
\[
    x = 0 \quad \text{ou} \quad x = -1
\]

    \item $7x^2 + x + 2 = 0$

    Discriminant :
\[
    \Delta = 1^2 - 4 \times 7 \times 2 = 1 - 56 = -55
\]
    Pas de solution réelle.

    \item $(x - 5)(x + 7) = 0$

    Solutions :
\[
    x = 5 \quad \text{ou} \quad x = -7
\]

    \item $(x - 3)(x + 1) = 6$

    Développons :
\[
    x^2 - 3x + x - 3 = 6 \quad \Rightarrow \quad x^2 - 2x - 9 = 0
\]
    Discriminant :
\[
    \Delta = (-2)^2 - 4 \times 1 \times (-9) = 4 + 36 = 40
\]
    Racines :
\[
    x = \dfrac{2 \pm 2\sqrt{10}}{2} = 1 \pm \sqrt{10}
\]

    \item $x^2 = \sqrt{2}$

    Solutions :
\[
    x = \pm (\sqrt{2})^{1/2} = \pm (\sqrt{2})^{1/2} = \pm 2^{1/4}
\]
    Simplifions :
\[
    x = \pm 2^{1/2} \times 2^{-1/2} = \pm 2^{(1 - 1)/2} = \pm 1
\]
    Mais cela n'est pas correct. En fait :
\[
    x = \pm \sqrt{\sqrt{2}} = \pm (\sqrt{2})^{1/2}
\]
    Donc :
\[
    x = \pm 2^{1/4}
\]
    \textbf{Réponse :} $x = \pm 2^{1/4}$
\end{enumerate}
\end{solutionbox}

\end{document}
